\chapter{Experimental Investigation of Mean Bias and Low-Bit Softmax Stability}
\label{chap:fp4_softmax_report}

\section{Problem Statement}

We investigate the interaction between activation anisotropy, mean bias, and
low-bit (FP4) quantization in Transformer attention.

In modern low-bit training pipelines, linear layers are computed in low precision
(FP4 or INT4), but nonlinear operators such as softmax are typically executed in
higher precision (e.g., BF16/FP16). This avoids catastrophic instability caused by
quantizing the probability simplex.

The central question of this investigation is:

\begin{center}
\emph{Is mean bias the dominant structural cause of softmax instability under FP4 quantization?}
\end{center}

To answer this question, we conducted a sequence of controlled experiments on
real model activations and attention states. The goal was to isolate:

\begin{enumerate}
\item The impact of mean bias on FP4 softmax error.
\item The relative effect of input quantization versus output quantization.
\item Whether mean removal alone stabilizes softmax in FP4.
\item Whether alternative attention operators improve robustness.
\end{enumerate}

All experiments were conducted using actual model tensors:
\[
Q, K \in \mathbb{R}^{1 \times H_{head} \times L \times d},
\]
with $L=1024$ tokens and $d=128$ per head.

FP4 quantization was simulated using blockwise scaling with percentile-based scale
selection and stochastic rounding.

\section{Experiment 1: FP4 Quantization of Activation Tensors}

\subsection{Objective}

To quantify whether removing mean bias improves softmax robustness when
queries and keys are quantized to FP4 but softmax output remains in FP16.

\subsection{Setup}

Two baselines were constructed:

\begin{itemize}
\item \textbf{Biased baseline:} raw $Q,K$ tensors.
\item \textbf{Mean-free baseline:} $Q,K$ centered along feature dimension.
\end{itemize}

For each baseline:

\begin{enumerate}
\item Compute reference softmax in FP16.
\item Quantize $Q,K$ to FP4.
\item Compute softmax in FP16.
\item Compare output against FP16 reference.
\end{enumerate}

Error metrics:

\[
\text{KL}, \quad L_1, \quad L_2, \quad \max |\Delta A|, \quad \text{tail log error}.
\]

\subsection{Results}

Representative results (Layer 3, Step 10k):

\begin{align*}
\textbf{Biased FP4 QK:} \quad & \text{KL}=3.81\times 10^{-2} \\
\textbf{Mean-free FP4 QK:} \quad & \text{KL}=3.32\times 10^{-2}
\end{align*}

Mean removal reduces KL divergence by approximately $13\%$,
and similarly reduces $L_1$, $L_2$, and tail error.

\subsection{Conclusion}

Mean removal improves softmax robustness under FP4 input quantization,
but the improvement is moderate. Instability does not vanish.

This indicates mean bias contributes to scale inflation,
but it is not the only mechanism of instability.

\section{Experiment 2: 2$\times$2 Softmax Quantization Study}

\subsection{Objective}

To separate the effect of input quantization (Q,K)
from output quantization (softmax probabilities).

We constructed a 2$\times$2 experimental grid:

\[
\text{Input: Biased vs Mean-free}
\quad \times \quad
\text{Output: FP16 vs FP4}.
\]

\subsection{Setup}

For each case:

\begin{enumerate}
\item Compute FP16 reference.
\item Quantize Q,K to FP4.
\item Compute attention.
\item Optionally quantize output probabilities to FP4.
\end{enumerate}

\subsection{Key Findings}

When softmax output is quantized directly to FP4
(using naïve probability quantization):

\[
\text{KL} \approx 22,
\quad
\text{out\_zero\_frac} \approx 0.999.
\]

This indicates catastrophic simplex collapse:
almost all attention mass is rounded to zero.

Using a log-ratio FP4 codec significantly reduces collapse:

\[
\text{KL} \approx 0.06,
\quad
\text{out\_zero\_frac} \approx 0.03.
\]

Mean removal improves KL by roughly $10$--$15\%$ in this regime.

\subsection{Conclusion}

The dominant instability under FP4 arises from
softmax output quantization, not merely input mean bias.

Mean bias affects scale selection,
but the primary structural problem is the difficulty of
representing probability simplices in FP4.

\section{Experiment 3: Laplacian Attention Operator}

\subsection{Objective}

To test whether replacing softmax attention with
a mean-neutral Laplacian attention update improves robustness.

Laplacian update:

\[
L = AV - V.
\]

\subsection{Findings}

When quantizing Q,K only:

\[
\|\Delta O\|_2 = \|\Delta L\|_2.
\]

When quantizing V as well:

Laplacian attention performs worse than standard attention,
because it exposes V-quantization error directly.

\subsection{Conclusion}

Difference-style updates do not mitigate FP4 instability.
They do not alter Jacobian sensitivity to A-noise.

\section{Experiment 4: Theoretical Analysis of Softmax Amplification}

We derived the three-stage instability mechanism:

\begin{center}
\textbf{Scale Hijacking} $\rightarrow$
\textbf{Variance Inflation} $\rightarrow$
\textbf{Dense Coupling Amplification}.
\end{center}

\subsection{Stage I: Scale Hijacking}

If logits decompose as

\[
z = m + r,
\]

and $\|m\|_\infty \gg \|r\|_\infty$,
then block scale satisfies

\[
s \approx \|m\|_\infty.
\]

Thus FP4 step size is governed by coherent mean bias.

\subsection{Stage II: Variance Inflation}

Stochastic rounding variance obeys

\[
\mathrm{Var}(\varepsilon) \le \Delta^2/4,
\quad \Delta \propto s.
\]

Hence

\[
\mathrm{Var}(\varepsilon) \propto \|m\|_\infty^2.
\]

\subsection{Stage III: Softmax Amplification}

Softmax Jacobian:

\[
J = \frac{1}{\tau}(\mathrm{diag}(y) - yy^\top).
\]

Error covariance propagates as

\[
\mathrm{Cov}(\Delta y) \approx J \Sigma_\varepsilon J^\top.
\]

In peaky regimes, variance from the dominant coordinate
is broadcast to all others.

\subsection{Hessian Drift}

Second-order Taylor expansion shows:

\[
\mathbb{E}[\hat y - y]
\propto \Sigma_\varepsilon,
\]

so quantization noise induces systematic bias.

\section{Global Conclusions}

From all experiments and derivations:

\begin{enumerate}
\item Mean bias contributes significantly to FP4 instability,
      especially through scale hijacking.
\item However, even mean-free inputs remain sensitive to FP4 quantization.
\item The dominant structural bottleneck is representing
      softmax probabilities in FP4.
\item Softmax amplification arises from dense coupling in its Jacobian.
\item Log-domain encoding dramatically improves stability.
\end{enumerate}

\section{Implications}

The instability of FP4 attention arises from a structural interaction:

\begin{equation}
\text{High-dimensional bias}
\;\times\;
\text{Blockwise scaling}
\;\times\;
\text{Simplex representation}
\;\times\;
\text{Softmax coupling}.
\end{equation}

Mean removal mitigates the first factor.
Improved probability encoding mitigates the third.
Jacobian-aware operator redesign is required to address the fourth.

\bigskip

These findings clarify that mean bias is a first-order contributor,
but full FP4 attention requires coordinated redesign of
scale selection, probability encoding, and operator structure.